\par\SourceLeaf{493}%
Inzwischen habe ich auf mehreren Wegen versucht, Nachrichten einzuholen, die in mehrfacher Hinsicht für die hannoversche Gradmessung wichtig sein werden. Diese gewinnt dadurch an Bedeutung, dass sie einerseits die ausgedehnte dänische Messung fortsetzt und andererseits noch viel weiter nach Süden fortgesetzt werden kann. Es freut mich daher, jetzt berichten zu können, dass ein beträchtlicher Teil der letztgenannten Arbeiten, die früher nur als möglich und wünschenswert dargestellt werden konnten, bereits tatsächlich ausgeführt worden ist. Die königlich preußische Regierung beabsichtigt nämlich, mehrere Hauptprovinzen der preußischen Monarchie vermessen zu lassen. Nach den heute als allein zulässig anerkannten Grundsätzen legt sie dieser Vermessung zunächst eine große Triangulation zugrunde. Sie hat bereits die Triangulation \textit{wirklich} ausführen lassen, die die preußischen Rheinprovinzen mit den älteren Provinzen verbindet und sich durch Nassau, Kurhessen, Thüringen und so weiter bis Halle erstreckt. Der königlich preußische Generalleutnant von \textsc{Müffling} hat mir diese Triangulation vollständig mitgeteilt. Nach eigener sorgfältiger Prüfung habe ich mich überzeugt, dass sie mit so großer Sorgfalt und Genauigkeit ausgeführt wurde, dass sie den Triangulationen bei den besten Gradmessungen an die Seite gestellt werden kann und einige sogar noch übertrifft. Diese Dreiecke führen zum Teil ganz nahe an den südlichen Grenzen des Königreichs Hannover vorbei. Wenn meine eigenen künftigen Messungen also gehörig daran angeschlossen werden, wozu mir Generalleutnant von \textsc{Müffling} bereits die nötigsten Auskünfte gegeben hat, wird die dänisch-hannoversche Gradmessung schon dadurch bis zur Sternwarte Seeberg bei Gotha und bis zum Inselsberg fortgesetzt sein. Nach Südwesten wird dadurch und durch die Darmstädter und französischen Messungen auch die Verbindung mit den Sternwarten von Mannheim und Paris erreicht sein.

Auch über die Messungen, die der französische Oberst \textsc{Epailly} während der französischen Besetzung in den Jahren 1803—1805 im Kurfürstentum Hannover ausführte, habe ich auf mehreren Wegen und zum Teil unmittelbar vom Dépot de la guerre mehrere wichtige, wenn auch unvollständige Nachrichten erhalten. Ein früherer Versuch, sie auf diplomatischem Weg vollständig mitgeteilt zu bekommen, war vergeblich. Ich habe noch einige Hoffnung, dass mir \textit{sämtliche} Dreiecke mitgeteilt werden. Sollte künftig einmal eine allgemeine Triangulation des ganzen Königreichs Hannover beschlossen werden, wofür die Gradmessung von Hamburg bis zur südlichen Grenze des Königreichs die solideste Grundlage liefern wird, so könnte der Besitz jener Messungen von \textsc{Epailly} unstreitig außerordentlich nützlich werden. Was mir davon mitgeteilt wurde, setzt mich zumindest in die Lage, eine solche mögliche Nutzung bei der Anordnung meiner eigenen Dreiecke im Voraus zu berücksichtigen und sie gewissermaßen vorzubereiten. — Unter den mir zugegangenen Aktenstücken über die Messung von \textsc{Epailly} befindet sich übrigens auch ein Bericht des Obersten \textsc{Epailly} selbst an General \textsc{Samson}. Dadurch habe ich freilich schon im Voraus von den großen \textit{Schwierigkeiten} der Triangulationen erfahren, die durch mancherlei örtliche Umstände besonders im Lüneburger Land entstehen. Diese Schwierigkeiten fand der französische Geodät selbst im ‘pays conquis’ fast unüberwindlich. Ich darf sie nur mit kräftiger Unterstützung der Regierung und der zuständigen Behörden zu überwinden hoffen.

Auch für mancherlei andere Erfordernisse der bevorstehenden Messungen habe ich im Lauf dieses Winters geeignete Vorkehrungen getroffen. Der erwartete große Theodolit von \textsc{Reichenbach} ist für die endgültige Messung der Winkel mit höchster Genauigkeit bestimmt. Für das Aufsuchen der Standpunkte, die Erkundung der von jedem aus sichtbaren Punkte und die vorläufigen Messungen ist aber ein anderer Theodolit erforderlich, der sich leicht transportieren, schnell aufstellen und einfach handhaben lässt. Dazu kommt, dass das Hauptinstrument größter Schonung bedarf. Zu diesem Zweck
